\subsection{A precise positive rule for the Schur-projector gate}
\begin{finding}[\fid{F23}: signed orbits suffice, but need not exist in pure Satake spaces]\label{f:signedorbits}
Let $M$ be a characteristic-zero $S_d$-module with a specified basis
$B$ permuted up to signs. For a label orbit $O$, choose $b_O$, let
$H_O$ be its label stabilizer, and define
$\chi_O(h)$ by $hb_O=\chi_O(h)b_O$. Then
\[
 M\simeq\bigoplus_O\operatorname{Ind}_{H_O}^{S_d}\chi_O,
\]
and
\[
 \dim M^{S_d}=\#\{O:\chi_O=1\},\qquad
 \dim\Hom_{S_d}(\mathrm{sgn},M)
 =\#\{O:\chi_O=\mathrm{sgn}|_{H_O}\}.
\]
For a primitive Young idempotent $e_\nu$,
\begin{equation}\label{eq:orbitprojector}
 \dim e_\nu M
 =\sum_O\dim\Hom_{H_O}(\Res^{S_d}_{H_O}[\nu],\chi_O).
\end{equation}
If every $H_O$ is a Young subgroup $\prod_i S_{a_i}$ and $\chi_O$ is a
product of trivial or sign characters on its blocks, this last sum
counts chains from $\varnothing$ to $\nu$: at block $i$ add a horizontal
$a_i$-strip for the trivial character, or a vertical $a_i$-strip for sign.

These are actual positive formulas for the corresponding Schur
multiplicity whenever the specified signed-basis data exist on
\[
 M_{\lambda,\mu,d}
 =\Hom_{\GL(V)}(S_\lambda V,(S_\mu V)^{\otimes d}).
\]
They do not always exist: for $\GL_2$, $\mu=(1)$, $d=3$ and
$\lambda=(2,1)$ this multiplicity module admits no signed-permutation
basis, although its natural geometric realization is
$\overline\Q_\ell(-1)^2$ with scalar Frobenius.
\end{finding}
\begin{proof}
The sign cocycle implies that $\chi_O:H_O\to\{\pm1\}$ is a character.
The map $k[S_d]\otimes_{k[H_O]}k_{\chi_O}\to kO$ sending
$g\otimes1$ to $gb_O$ is well defined, equivariant and surjective.
Both spaces have dimension $[S_d:H_O]$, so it is an isomorphism.
The distinct orbit spans are a direct sum. Frobenius reciprocity gives
\eqref{eq:orbitprojector}, and the trivial/sign cases have a
one-dimensional Hom precisely when the indicated characters agree.
In $M=\bigoplus_\rho[\rho]\otimes W_\rho$, a primitive idempotent has
rank one on $[\nu]$, hence image dimension $\dim W_\nu$.
A central idempotent would have rank $\dim[\nu]\dim W_\nu$ instead.
When the stabilizer is the stated Young subgroup, the Frobenius
characteristic is a product of $h_{a_i}$ and $e_{a_i}$. The classical
Pieri rule \cite[Theorem~7.10, p.~16]{Pieri} expands this product by exactly
the horizontal/vertical chains just described. This proves the positive
chain formula, without presupposing the desired projector-image basis.

For the counterexample, Schur--Weyl duality gives
\[
 (\overline\Q_\ell^2)^{\otimes3}
 =S_{(3)}V\otimes[(3)]\oplus S_{(2,1)}V\otimes[(2,1)].
\]
Thus the indicated multiplicity is the faithful two-dimensional
standard representation of $S_3$. It is faithful because it is the
sum-zero subspace of the permutation representation on three letters,
whose complementary all-ones line is fixed. A signed-permutation basis
would embed $S_3$ into the group of $2\times2$ signed permutation
matrices, of order $2^2\cdot2!=8$. Lagrange's theorem excludes this.

Here is the geometric identification in full. The threefold minuscule
convolution variety for $\GL_2$ is an iterated $\PP^1$-bundle, smooth of
dimension three. Over the lattice
$t^2\F_q[[t]]e_1\oplus t\F_q[[t]]e_2$, the quotient has dimension
three and multiplication by $t$ is nilpotent of type $(2,1)$. The fiber
parametrizes its full stable flags $L\subset P$. Write
$I=\operatorname{im}N$, $K=\ker N$, of dimensions one and two.
Stability requires $L\subset K$, $I\subset P$, and $N(P)\subset L$.
If $P\ne K$, then $N(P)=I$ and $L=I$; if $P=K$, every line
$L\subset K$ works. Thus the reduced fiber is
\[
 \{L=I,\ P\supset I\}\ \cup\ \{P=K,\ L\subset K\}
 =\PP^1\cup_{\mathrm{pt}}\PP^1.
\]
Both components and the intersection are rational over $\F_q$.
Its $H^2$ is therefore $\overline\Q_\ell(-1)^2$. Since the source
dimension is three and $D_{(2,1)}=1$, proper base change identifies this
group with the degree $-1$ stalk of the pure perverse convolution.
The intermediate extension $\IC_{(3)}$ has no degree
$-\dim\Gr_{(2,1)}=-1$ stalk on this boundary stratum, by its strict
support condition. This stalk is exactly the $\IC_{(2,1)}$ multiplicity
module. Corrected Satake symmetry gives the standard $S_3$ action just
used, while Frobenius is scalar $q$. This proves every assertion about
the example.
\end{proof}
The obstruction is to signs $\{\pm1\}$, not arbitrary scalar-monomial
bases: the standard $S_3$ representation is induced from a nontrivial
character of its cyclic subgroup of order three. Pure Tate Frobenius
does not restrict the phases of this separate commuting group action.
The example also has classical positive tableau formulas; it refutes
only the asserted universal signed-orbit route.

\paragraph{Exact encoding conditions for a counting rule.}
For inputs of binary length $L$, suppose orbit representatives have
unique polynomial-length encodings with polynomial-time validity tests.
For the trivial/sign formula, suppose a polynomial-size stabilizer
generating set and its character values are computable in polynomial
time, including the cost of permutation/sign encoding when $d$ is
binary. The character test on generators makes the orbit label itself
a counted witness. For the mixed-Young formula, additionally require
polynomially many blocks and rows, with polynomial-bit block sizes and
parts, and a polynomial-time procedure computing the block sizes and
trivial/sign types from the orbit label. A witness consists of that label
and its chain of partitions.
All parts are bounded by $d$, so the witness is polynomial length.
For $\alpha\subset\beta$, a horizontal strip is checked by
$\beta_i\ge\alpha_i\ge\beta_{i+1}$; a vertical strip by
$\beta_i-\alpha_i\in\{0,1\}$. Sizes and partitions are checked by
integer arithmetic. These tests count exactly the displayed chains,
giving a conditional \(\#P\) rule. They do not provide a single
polytope, an Ehrhart family or PH1's oracle from purity alone.

\paragraph{A known binary family, with its exact small polytope.}
For $\dim V=2$, Hermite reciprocity and Clebsch--Gordan with swap parity
give the classical formula
\[
 \Sym^d(\Sym^2 V)\simeq\Sym^2(\Sym^dV)
       \simeq\bigoplus_{j=0}^{\lfloor d/2\rfloor}S_{(2d-2j,2j)}V.
\]
The exact sources are Goodman--Wallach, Theorem~4.1.18, p.~191,
and Donley--Kim, Proposition~3.1, equation~(3.4), p.~3, and
Proposition~8.1, equation~(8.2), p.~11 \cite{GW,DK}.
These known theorems are not re-proved here. The passage from $SL_2$
to $GL_2$ fixes total polynomial degree $2d$; in the summand
$S_{(2d-r,r)}$, a highest-weight vector is
$(x_1y_2-x_2y_1)^r x_1^{d-r}x_2^{d-r}$ and swap acts by $(-1)^r$.
For $\lambda=(2d-r,r)$, $0\le r\le d$, the specialized consequence is
\[
 a_\lambda(d[2])=\mathbf1_{2\mid r},\qquad
 a_{N\lambda}((Nd)[2])
 =\#(N P_{d,r}\cap\Z),\quad P_{d,r}=\{r/2\}.
\]
The point is given by $2x\le r$, $-2x\le-r$, so has an explicit
polynomial-bit description and separation test. This is a restricted
fixed-inner-partition Ehrhart formula, not general PH1. On odd-$r$
rays its ordinary support is unsaturated, but the zero odd residue
polynomial is permitted by VI Definition~1.2.4, so the VI saturation
index remains zero, exactly as distinguished in \fid{F18}.

